\lecture{17}{3. November 2025}{Fluid Flows About Immersed Bodies -- Drag, Lift and Airfoils}

\begin{problem}{9.28}
  The lift and drag coefficients for a rectangular wing with a \qty{10}{\m} chord at takeoff are \num{1,0} and \num{0,05}, respectively.
  Determine the span needed to lift \qty{3560}{\kN} at a take-off speed of \qty{282}{\km\per\hour}.
  Determine the wing drag at this condition.
\end{problem}

\begin{assumptions}
  \item The temperature is \qty{20}{\celsius} 
\end{assumptions}

\begin{derivation}
  The general formula for the lift force is
  \begin{equation*}
    F_L = \frac{1}{2} \rho C_L V^2 A_p
  \end{equation*}
  where $A_p$ is the planform area given by
  \begin{equation*}
    A_p = s \cdot l
  \end{equation*}
  where $s$ is span and $l$ is the chord length.
  Substituting this into the formula for the lift and solving for the required span $s$ we obtain
  \begin{align*}
    F_L &= \frac{1}{2} \rho C_L V^2 s \cdot l \\
    s &= \frac{2F_L}{\rho C_L V^2 \cdot l}
  .\end{align*}
  For air at \qty{20}{\celsius} the density can, in Table A.10 in appendix A of the course literature be found to be approximately $\rho_{\text{air}} = \qty{1,21}{\kg\per\cubic\m}$.
  Plugging in the numbers we get
  \begin{equation*}
    s = \frac{2 \cdot \qty{3560}{kN}}{\qty{1,21}{\kg\per\cubic\m} \cdot \num{1,0} \cdot \left( \qty{282}{\km\per\hour}  \right)^2 \cdot \qty{10}{\m} } = \qty{95,9}{\m}
  .\end{equation*}
  The drag force is given as
  \begin{equation*}
    F_D = \frac{1}{2}\rho C_D A U^2
  .\end{equation*}
  Plugging in the numbers we get
  \begin{equation*}
    F_D = \frac{1}{2} \cdot \qty{1,21}{\kg\per\cubic\m} \cdot \num{0,05} \cdot \qty{95,9}{\m} \cdot \qty{10}{\m} \cdot \left( \qty{282}{\km\per\hour}  \right)^2 = \qty{178}{\kN} 
  .\end{equation*}
\end{derivation}

\finalresult{
\begin{aligned}
  s &= \qty{95,9}{\m}  \\
  F_{D} &= \qty{178}{\kN} 
.\end{aligned}
}

\begin{problem}{9.30}
  The mean velocity over the top of a wing with a \qty{1,8}{\m} chord moving through air at \qty{33,5}{\m\per\s} is \qty{40}{\m\per\s}, and that over the bottom of the wing is \qty{31}{\m\per\s}.
  Determine the lift per meter of span and the lift coefficient.
\end{problem}

\begin{assumptions}
  \item Steady flow at \qty{20}{\celsius} 
\end{assumptions}

\begin{derivation}
  We will use the Bernoulli equation to calculate the difference in pressure over the top of the wing as
  \begin{align*}
    p_t + \frac{1}{2} \rho U_t^2 &= p_a + \frac{1}{2} \rho U_a^2 \\
    p_t - p_a &= \frac{1}{2} \rho \left( U_a^2 - U_t^2 \right) \\
              &= \frac{1}{2} \cdot \qty{1,225}{\kg\per\cubic\m} \cdot \left( \left( \qty{33,5}{\m\per\s} \right)^2 - \left( \qty{40}{\m\per\s}  \right)^2 \right) \\
              &= \qty{-293}{\Pa} 
  .\end{align*}
  Doing the same over the bottom of the wing we get
  \begin{align*}
    p_b - p_a &= \frac{1}{2} \rho \left( U_a^2 - U_b^2 \right) \\
    &= \frac{1}{2} \cdot \qty{1,225}{\kg\per\cubic\m} \cdot \left( \left( \qty{33,5}{\m\per\s}  \right)^2 - \left( \qty{31}{\m\per\s}  \right)^2 \right) \\
    &= \qty{98,77}{\Pa} 
  .\end{align*}
  The difference in pressure between the top and bottom is then
  \begin{equation*}
    p_b - p_t = \qty{98,77}{\Pa} - \left( \qty{-293}{\Pa}  \right) = \qty{391,77}{\Pa} 
  .\end{equation*}
  The lift force is the pressure difference times the area, hence the lift per meter of spa, where $l$ is the length of the airfoil chord is
  \begin{equation*}
    F_L = \left( p_b - p_t \right) l = \qty{391,77}{\Pa} \cdot \qty{1,8}{\m} = \qty{705,186}{\N\per\m} 
  .\end{equation*}
  The lift coefficient can then be determined as
  \begin{equation*}
    C_L = \frac{F_L}{\frac{1}{2} \cdot \rho \cdot V^2 \cdot l} = \frac{\qty{705,186}{\N\per\m} }{\frac{1}{2}\cdot \qty{1,225}{\kg\per\cubic\m} \cdot \left( \qty{33,5}{\m\per\s}  \right)^2 \cdot \qty{1,8}{\m} } = \num{0,567} 
  .\end{equation*}
  
\end{derivation}

\finalresult{
  \begin{aligned}
    F_L &= \qty{705,186}{\N\per\m} \\
    C_L &= \num{0,567} 
  .\end{aligned}
}


\begin{problem}{9.6RP}
  A high voltage transmission line cable \qty{10}{\cm} in diameter runs between two towers \qty{500}{\m} apart.
  Determine the force on the cable for a wind speed of \qty{15}{\m\per\s}.
  Determine whether the total force would be less of four \qty{5}{\cm} cables with the same electrical capacity were used.
\end{problem}

\begin{assumptions}
  \item Steady flow at \qty{20}{\celsius} 
  \item The cable is a rigid circular cylinder
\end{assumptions}

\begin{derivation}
  The drag coefficient is 
  \begin{equation*}
    C_D = \frac{F_D}{\frac{1}{2} \rho U^2 A}
  .\end{equation*}
  As the drag coefficient of a circular cylinder is a function of the Reynolds number we must first determine the Reynolds number as:
  \begin{equation*}
    \mathrm{Re} = \frac{\rho \cdot U \cdot D}{\mu} = \frac{\qty{1,21}{\kg\per\cubic\m} \cdot \qty{15}{\m\per\s} \cdot \qty{10}{\cm}}{\qty{1,81e-5}{\N\s\per\square\m}} = \num{e5} 
  .\end{equation*}
  From Fig. 9.13 of the course literature, the drag coefficient is determined to be about $C_D = \num{1,0}$ for a Reynolds number of $\mathrm{Re} = \num{e5}$.
  As the frontal area of the cylinder-shaped cable is $A = DL$ the drag force is
  \begin{equation*}
    F_D = \frac{1}{2} \rho C_D U^2 A = \frac{1}{2} \cdot \qty{1,21}{\kg\per\cubic\m} \cdot \left( \qty{15}{\m\per\s}  \right)^2 \cdot \qty{10}{\cm} \cdot \qty{500}{\m}  = \qty{6,8}{\kN} 
  .\end{equation*}
  For the four smaller cables the Reynolds number is:
  \begin{equation*}
    \mathrm{Re} = \frac{\qty{1,21}{\kg\per\cubic\m} \cdot \qty{15}{\m\per\s} \cdot \qty{5}{\cm} }{\qty{1,81e-5}{\N\s\per\square\m}} = \num{5e4} 
  .\end{equation*}
  This also corresponds to a drag coefficient of approximately $C_D = \num{1,0}$ per Fig. 9.13.
  Thus:
  \begin{equation*}
    F_D = \frac{1}{2} \rho C_D U^2 A = \frac{1}{2} \cdot \qty{1,21}{\kg\per\cubic\m} \cdot \left( \qty{15}{\m\per\s}  \right)^2 \cdot \qty{5}{\cm} \cdot \qty{500}{\m} = \qty{3,4}{\kN} 
  .\end{equation*}
  This is the drag force on each cable.
  For all four \qty{5}{\cm}  cables the total drag force is
  \begin{equation*}
    \left( F_{D} \right)_{\text{4 cable}} = 4 \cdot \qty{3,4}{\kN} = \qty{13,6}{\kN} 
  .\end{equation*}
\end{derivation}

\finalresult{
  \begin{aligned}
    F_D &= \qty{6,8}{\kN}  \\
    \left( F_{D} \right)_{\text{4 cable}} &= \qty{13,6}{\kN}
  .\end{aligned}
}
